\section{Travaux connexes}
\label{sec:related}

\textbf{Classification en endoscopie.}
Kvasir~\cite{pogorelov2017kvasir} et HyperKvasir~\cite{borgli2020hyperkvasir} sont des références pour la classification d'images digestives.
Un travail antérieur dans notre dépôt~\cite{assif2026endoscopy} compare ResNet-50, SE-ResNet-50 et fusion couleur--texture sur HyperKvasir (23 classes), avec une macro-F1 test $\approx$0,58--0,59 (B2--B6).
Nous ciblons une taxonomie \emph{à trois classes} alignée sur la pratique (repère, muqueuse, polype) plutôt que sur 23 labels académiques.

\textbf{Attention canal.}
SE-Net~\cite{hu2018senet} recalibre les canaux via une moyenne globale (GAP) puis excitation.
Sous masquage JEPA, les zéros de padding biaisent la GAP ; nous proposons un pooling pondéré par le masque (section~\ref{sec:maskse}).

\textbf{JEPA et auto-supervision.}
I-JEPA~\cite{assran2023ijepa} repose sur des Vision Transformers ; nous employons un \emph{CNN-JEPA} avec masquage par blocs et prédicteur convolutif sur cartes de features---plus adapté au déploiement léger qu'un foundation model ViT.
Nous ne prétendons pas inventer le pré-entraînement convolutif par prédiction latente ; notre originalité porte sur l'\textbf{intégration SE sensible au masque} et le \textbf{protocole de benchmark reproductible} pour la classification coloscopique.
Des travaux proches incluent BYOL~\cite{grill2020byol}.

\textbf{Détection de polypes.}
Les systèmes CAD~\cite{tajbakhsh2016cad} mettent l'accent sur la \emph{sensibilité} aux polypes.
Nous rapportons le rappel polype en plus de la macro-F1.
La classe \emph{polype} est unique (Kvasir + HyperKvasir \texttt{polyps} et \texttt{dyed-lifted-polyps}) ; le typage histologique est reporté aux travaux futurs.
